\documentclass[11pt,a4paper]{article}
\usepackage[utf8]{inputenc}
\usepackage{amsmath,amssymb}
\usepackage[margin=1in]{geometry}
\usepackage{charter}
\usepackage{microtype}
\usepackage{hyperref}
\usepackage{xcolor}

\hypersetup{
    colorlinks=true,
    linkcolor=blue!70!black,
    urlcolor=blue!70!black
}

\setlength{\parindent}{0pt}
\setlength{\parskip}{0.7em}

\begin{document}

\begin{flushright}
\textbf{Date:} August 31, 2026
\end{flushright}

\textbf{Prof. Kenneth M. Merz, Jr.}\\
Editor-in-Chief\\
\textit{Journal of Chemical Information and Modeling}\\
American Chemical Society

\bigskip

\textbf{Subject: Resubmission of Manuscript ci-2026-02951g to \textit{Journal of Chemical Information and Modeling}}

\bigskip

Dear Prof. Merz and Editorial Board Members,

We are pleased to resubmit our manuscript entitled \textbf{``Physics-Gated Interpretable Machine Learning and Symbolic Regression for Double Perovskites: Overcoming Theoretical Limits Beyond 0D Compositional Models''} (Manuscript ID: \textbf{ci-2026-02951g}) for consideration as a full Research Article in the \textit{Journal of Chemical Information and Modeling} (JCIM).

\textbf{Response to Editorial Formatting \& Resubmission Notes:}
In accordance with the editorial unsubmission notice for administrative technical checks, we have completed all required formatting updates:
\begin{enumerate}
    \item \textbf{Standalone Supporting Information for Publication}: A dedicated, standalone PDF (\texttt{supporting\_information.pdf}) has been prepared with full matching manuscript title, authors, institutional affiliations, corresponding author contact emails on Page S1, and standard sequential $S$-numbering (Pages S1--S9, Tables S1--S5, Figures S1--S4).
    \item \textbf{Continuous Narrative Introduction}: The main manuscript Introduction has been formatted into a continuous academic narrative without subheadings or intrusive tables per ACS style.
    \item \textbf{Complete Submission Package}: The main manuscript, supporting information, and table of contents graphic have been updated and designated per ACS Publishing Center guidelines.
\end{enumerate}

\textbf{Significance and Core Advances of the Work:}
Predicting the ground-state properties of complex multinary materials like double perovskites ($A_2BB'O_6$) using purely compositional (0D) descriptors is a major challenge in chemical informatics. While 3D graph neural networks (CGNNs) offer high fidelity, they require computationally expensive, pre-relaxed DFT coordinates and function as uninterpretable black-boxes. Conversely, foundational 0D symbolic models face theoretical accuracy ceilings ($R^2_{\text{limit}} \approx 65\%$ for formation energy $\Delta E_f$, $60\%$ for magnetization $M$, $50\%$ for band gap $E_g$, and $25\%$ for energy above hull $E_{\text{hull}}$) and collapse on zero-inflation point-mass densities.

In this study, we introduce a novel \textbf{Physics-Gated Interpretable Machine Learning Architecture} that breaks these theoretical literature ceilings using 100\% 0D descriptors. Our key contributions include:
\begin{itemize}
    \item \textbf{Physics-Gated Domain Routers}: Resolving zero-inflation and derivative discontinuities at metallic/insulating and magnetic phase boundaries via high-margin hurdle models ($C=200.0$) and soft-sigmoidal gating functions.
    \item \textbf{Quantum \& Thermodynamic Engines}: Deriving closed-form Harrison tight-binding inter-atomic matrix elements and single-perovskite tie-line proxies ($D_{\text{hull\_proxy}}$) to boost 0D accuracy without leaking 3D coordinate data.
    \item \textbf{Surpassing SOTA Literature Limits}: Achieving \textbf{109.62\%} of the theoretical limit for formation energy ($\Delta E_f\ R^2 = 71.26\%$), \textbf{103.72\%} for total magnetization ($M\ R^2 = 62.23\%$), and \textbf{101.42\%} for band gap ($E_g\ R^2 = 50.71\%$).
    \item \textbf{Rigorous Multi-Seed Validation}: Validated across 25 independent random seeds and out-of-distribution splits on 2,000 and 5,000 double perovskite compositions.
    \item \textbf{Closed-Form Analytical Equations}: Discovering human-interpretable symbolic equations directly applicable for prospective materials screening.
\end{itemize}

\textbf{Fit for the \textit{Journal of Chemical Information and Modeling}:}
This work directly aligns with JCIM's core scope at the intersection of chemical informatics, data-driven materials modeling, and interpretable physical machine learning.

\textbf{Declarations:}
\begin{itemize}
    \item This manuscript is original, has not been published previously, and is not under consideration elsewhere.
    \item All authors have approved the manuscript and agree to its resubmission to JCIM.
    \item All datasets and Python source code are openly available on GitHub at \url{https://github.com/nihal-gazi/double_perovskite_compositional_interpretable}.
    \item The authors declare no competing financial interests.
\end{itemize}

\textbf{Suggested Independent Reviewers:}
\begin{enumerate}
    \item \textbf{Prof. Aron Walsh} (Imperial College London, UK), \textit{Email:} \url{a.walsh@imperial.ac.uk} -- World-renowned expert in perovskite electronic structure, crystal chemistry, and tight-binding physical modeling.
    \item \textbf{Dr. Ghanshyam Pilania} (Los Alamos National Laboratory, USA), \textit{Email:} \url{gpilania@lanl.gov} -- Pioneer in compositional (0D) feature engineering, band gap machine learning, and physics-informed materials informatics for perovskites.
    \item \textbf{Prof. Bryce Meredig} (Stanford University / Citrine Informatics, USA), \textit{Email:} \url{bryce@citrine.io} -- Leading authority in domain-informed compositional descriptors and interpretable materials screening over black-box networks.
    \item \textbf{Prof. Rampi Ramprasad} (Georgia Institute of Technology, USA), \textit{Email:} \url{rampi.ramprasad@mse.gatech.edu} -- Distinguished expert in machine learning for functional materials design, physical property mapping, and chemical informatics.
\end{enumerate}

Thank you very much for your time, consideration, and continued handling of our manuscript.

\bigskip

Sincerely yours,

\textbf{Dr. Soumyadipta Pal} (Corresponding Author) \\
University of Engineering and Management (UEM), Kolkata, India; Email: \url{soumyadipta.pal@gmail.com}

\textbf{Dr. Subarna Datta} (Corresponding Author) \\
Institute of Engineering \& Management (IEM), Kolkata, India; Email: \url{subarna.iem87@gmail.com}

\end{document}
